\section{四次のKoszul検査を通る非KoszulなFrobenius境界}
\label{sec:nonkoszul-boundary}
有限Frobenius性と元代数のKoszul性を分ける必要は形式上だけではない。
以下では対称チャートの閉包内で、この差を五次に初めて検出する具体例を与える。

\begin{theorem}[有限Frobenius領域内の非対称化境界]\label{thm:nonkoszul-symmetric-boundary}
$Z_{\mathrm{sym}}$には有限Frobenius候補$[w_0]$であって対称化も交代化もできない点が存在する。
その第三階数は6であり、元の四窓代数$C$は全始点で四次Koszul検査を通るが、Koszulではない。
特に対称チャートの飽和像$S$は、有限Frobenius低階数候補領域の中で相対的に閉ではない。
ただし$C$は本研究の標準型Koszul AS正則対象ではなく、未分類AS枝の存在証明にはならない。
\end{theorem}
\begin{proof}
可換二次式の四次元部分空間
\[
 L_t=\langle e_0^2-e_0e_3,\ e_1^2-e_1e_3,\ e_2^2-e_2e_3,
                  \ e_3(e_0+e_1+e_2+t e_3)\rangle
\]
を取る。最初の三式の射影共通零点は
$[\epsilon_0:\epsilon_1:\epsilon_2:1]$, $\epsilon_i\in\{0,1\}$の八点である。
最後の式の値は$\epsilon_0+\epsilon_1+\epsilon_2+t$なので、
$t\notin\{0,-1,-2,-3\}$なら$L_t$は基点自由な完全交叉である。
$t=0$では共通零点$[0:0:0:1]$が残る。

$B_t=\operatorname{Sym}(V^*)/(L_t)$と置く。
対称三テンソル$b=(b_{abc})$が$(B_t)_3^*$に属する条件は
\[
 b_{iia}-b_{i3a}=0\quad(0\leq i\leq2,\ 0\leq a\leq3),\qquad
 \sum_{i=0}^2 b_{i3a}+t b_{33a}=0\quad(0\leq a\leq3).
\]
20個の対称座標に対するこの16本の式に、
$012,133,233,333$の四座標を指定する式を加えた20次行列の行列式は$-2$である
（行順序は再現コードの通り）。
従って全ての$t$で$(B_t)_3$は4次元であり、その双対基底$b_0(t),\ldots,b_3(t)$を
上記四座標が標準基底となるよう一意に多項式で選べる。
因子$0,1,3$に三テンソルを、因子2に対応する基底を置いて
\[
 (w_t)_{abjc}=(b_j(t))_{abc}
\]
と定める。これは全$t$で非零な多項式テンソル族である。
完全交叉となる$t$では、その対称最上traceテンソルの一因子縮約は
$(B_t)_3^*$の基底を与えるので、因子2の可逆変換で$w_t$へ移る。
従って一般の$w_t$は$S$に属し、$w_0\in Z_{\mathrm{sym}}$である。

$t=0$の有理数線型計算は次の通りである：
\[
\begin{array}{c|c}
検査 & 値\\\hline
一因子階数 & (4,4,4,4)\\
三つの二対二階数 & (6,6,6)\\
\dim W_{i,3},\ \dim W_{i,4} & (4,1),(4,1),(4,1),(4,1)\\
\dim C_{i,i+4} & (35,35,35,35)\\
\dim C_{i,i+5} & (63,56,56,56)
\end{array}
\]
ここで三・四次交叉は復元した二窓の直交関係を挿入した行列の核で計算し、
最後の二行は元の六関係を挿入した行列の余核で計算する。
とくに五次挿入行列の階数は$(961,968,968,968)$である。
全係数・行列・Gauss消去の再現は
\path{computations/verify_nonkoszul_boundary.py}に置く。
テンソル復元定理により$w_0$は有限Frobenius候補であり、相補次数の対合も完全である。
さらに線型対称化空間は
\[
 S_1(w_0)=S_3(w_0)=kI,\qquad S_2(w_0)=kE_{33}
\]
となるので対称化できない。符号付き三空間の可逆性判定も失敗するので交代化もできない。
四次の値35は全ての$\delta_i$が零であることを意味する。
しかしKoszulなら、双対のHilbert列$(1,4,6,4,1)$から
$h_C(z)=(1-z)^{-4}$となり、全ての五次元数は56でなければならない。
始点0の値63がこれに反する。より正確には五次の$d_3$階数は全始点で36であり、
命題\ref{prop:quintic-koszul-obstruction}により$H_3=0$、$H_2$の次元は$(7,0,0,0)$である。
\end{proof}

\begin{remark}
この例は有限Frobenius性と四次Koszul切断の十分性をともに否定する。
五次での二階数条件の十分性や、一般の有限次数判定は主張しない。
また$S$の有限Frobenius領域内の非閉性と、Koszul AS領域内での非閉性を区別する。
後者、および対称チャート外の低階数AS対象の存在はなお未解決である。
\end{remark}
